\section{Homogeneity and a superlinear incidence estimate}
\label{sec:v80-incidence}
The inverse incidences in \eqref{eq:v79-optical-gradient} can be
summed on the complete source. The required improvement is not a
bound for multiplication by $c^{-1}$ on the strong space. It is a
stronger estimate for the real-frequency physical powers, followed
by a summable, width-dependent remainder.

\subsection{The precise spectral input}
Write $Q=\mathscr T_{R,\xi}$ for the uncentered four-coordinate
operator, with $\xi=(u_1,u_2,b,v)$ real. The densities in the
collision spaces are relative to $\nu$; the initial vector is
$\mathbf1$, and
\begin{equation}\label{eq:v80-physical-modulus}
 Q^j\mathbf1(x)=\exp\bigl(i\xi\cdot S_j f(T_R^{-j}x)\bigr),
                  \qquad |Q^j\mathbf1|=1
\end{equation}
on the regular physical realization. Here $S_jf$ denotes the full
record, not the roof alone. This equality would not justify a
cosine factor in the stable-curve norm.

We use the following source-pinned statements. The norm definitions
are Supplement equations~\eqref{supp-eq:v35-weak-norm}--\eqref{supp-eq:v35-unstable-norm};
the full-occupation Doeblin--Fortet inequality and fixed-band spectrum
are in Supplement Theorem~\ref{supp-thm:v40-polynomial-LY} and the
moving-peak construction used in Supplement
Lemma~\ref{supp-lem:v48-weak-interpolation}. The incidence multiplier
$M_s=M_{\one_{\{c<s\}}}$ satisfies
\begin{equation}\label{eq:v80-inherited-multiplier}
 \|M_s h\|\le C\|h\|,\qquad
 |M_s h|_w\le Cs^{1/8}\|h\|.
\end{equation}
This is Supplement Proposition~\ref{supp-prop:v49-physical-multiplier}.
On a fixed finite collection of peak charts one has
\begin{align}
 Q^j&=\lambda^j\Pi+N^j,&\|N^j\|&\le C\rho^j,&\|Q^j\|&\le C,
                                                        \notag\\
 \|Q^j h\|&\le C\sigma^j\|h\|+C(1+j)|h|_w,&
 |Q^j h|_w&\le C(1+j)|h|_w,                  \label{eq:v80-spectral-input}
\end{align}
where $\rho,\sigma<1$. The fixed-band complement has a uniform
spectral gap. The four-frequency Gaussian integral on each peak is
$O(m^{-2})$. These are inherited spectral theorems, not consequences
of the coarea argument.

After shrinking the neighborhoods of unit-modulus points, we may take
$|\lambda|\ge r_*$ on the retained charts with
\begin{equation}\label{eq:v80-peak-choice}
 \frac{\log(1/r_*)}{\log(1/\sigma)}<\frac1{32},\qquad
 \frac{\log(1/r_*)}{\log(1/\rho)}<\frac1{24}.
\end{equation}
The removed compact portions have a spectral gap and join the
complement. This is a fixed refinement, independent of $m,s$ and
of the eventual outer band. No arithmetic branch is discarded.

\begin{lemma}[Real powers on homogeneous grazing strips]
\label{lem:v80-homogeneous-weak}
For all $j\ge0$ and sufficiently small $s>0$, uniformly on the
retained real peak charts,
\begin{equation}\label{eq:v80-homogeneous-weak}
             |M_s Q^j\mathbf1|_w\le Cs^{3/2}.
\end{equation}
The exponent comes from the length of a homogeneous stable curve,
not from replacing the density convention by Lebesgue density.
\end{lemma}
\begin{proof}
The quadratic homogeneity strips have angular widths comparable to
$k^{-3}$ and incidences comparable to $k^{-2}$. In the standard
coordinates they lie between
$\pi/2-(k+1)^{-2}$ and $\pi/2-k^{-2}$, or their negative counterparts.
A stable graph has slope bounded away from zero and infinity.
Consequently a homogeneous stable curve in the $k$th strip has
length at most $Ck^{-3}$. If that strip meets $\{c<s\}$, then
$k\ge C^{-1}s^{-1/2}$; the same estimate includes a threshold
crossing the strip. Thus the entire tested intersection has length
at most $Cs^{3/2}$. For small $s$, the central strips do not meet it.
Equation~\eqref{eq:v80-physical-modulus} and the supremum bound of a
weak test now prove \eqref{eq:v80-homogeneous-weak}. This proof applies
to the physical branchwise realization; the bounded multiplier and
physical realization in \eqref{eq:v80-inherited-multiplier} give the
same element of the weak completion.
\end{proof}

\begin{lemma}[Interpolation close to the weak endpoint]
\label{lem:v80-interpolation}
On the fixed charts chosen above,
\begin{equation}\label{eq:v80-interpolation}
 \|\Pi h\|+\sup_{a\ge0}|\ell(Q^a h)|
                  \le C\|h\|^{1/8}|h|_w^{7/8}.
\end{equation}
Moreover
\begin{equation}\label{eq:v80-projected-incidence}
 |M_s\Pi\mathbf1|_w\le Cs^{23/16}.
\end{equation}
\end{lemma}
\begin{proof}
Normalize $\|h\|=1$ and put $t=|h|_w$. For $0<t<1/2$ choose
$L$ with $\sigma^L\le t<\sigma^{L-1}$. Commutation and
\eqref{eq:v80-spectral-input} give
\[
 \|\Pi h\|\le C r_*^{-L}\{\sigma^L+(1+L)t\}
             \le C(1+|\log t|)t^{31/32}\le Ct^{7/8}.
\]
For $a\le L$ the weak power bound controls the terminal pairing;
for $a>L$ apply the strong inequality at $L$ and power boundedness
afterwards. Both costs are at most $C(1+L)t\le Ct^{7/8}$.
The cases of $t$ bounded below and $t=0$ follow, respectively, by
power boundedness and by letting $L\to\infty$. Homogeneity gives
\eqref{eq:v80-interpolation}.

For \eqref{eq:v80-projected-incidence}, choose instead
$L=\lceil(3/2)\log(1/s)/\log(1/\rho)\rceil$. The exact identity
\[
 M_s\Pi\mathbf1
    =\lambda^{-L}M_sQ^L\mathbf1-\lambda^{-L}M_sN^L\mathbf1
\]
uses $N^L=Q^L-\lambda^L\Pi$. Its weak norm is at most
$Cr_*^{-L}(s^{3/2}+s^{1/8}\rho^L)$.
Here $\rho^L\le s^{3/2}$ and
$r_*^{-L}\le Cs^{-1/16}$ by \eqref{eq:v80-peak-choice}. This proves
the assertion. The strong norm of $M_s\Pi\mathbf1$ remains bounded.
This argument does not assign modulus one to the continued eigenvector.
\end{proof}

\begin{theorem}[A capped superlinear marked-incidence bound]
\label{thm:v80-superlinear-incidence}
There is $C$ such that, for every fixed auxiliary roof band $B_*\ge B_0$,
there are $C_{B_*}<\infty$ and $\rho_{B_*}<1$ satisfying
\begin{align}
 &m^2\nu\{K_m^{\rm c}=k,A_m=l,L_m\in J,\ c(T_R^jx)<s\}
                                                        \notag\\
 &\quad\le Cs^{5/4}(|J|+B_*^{-1})
  +C_{B_*}(1+B_*|J|)m^2
                   \min\{\rho_{B_*}^{m/2},s^2\}.       \label{eq:v80-incidence-local}
\end{align}
The estimate is uniform in $R$, $m\ge2$, $0\le j\le m$, exact
labels, intervals $J$ and $0<s<s_0$. In particular the mark and
threshold may vary with $m$. The leading constant is independent
of the previously fixed outer band.
\end{theorem}
\begin{proof}
The exact characteristic pairing is
$\ell(Q^{m-j}M_sQ^j\mathbf1)$. If the left block is long, expand
that block and apply \eqref{eq:v80-interpolation} and
\eqref{eq:v80-homogeneous-weak} to $M_sQ^j\mathbf1$.
The leading amplitude is at most $Cs^{21/16}$. If the right block
is long, its leading term is controlled by
\eqref{eq:v80-interpolation} and \eqref{eq:v80-projected-incidence};
its amplitude is at most $Cs^{161/128}$. Both are bounded by
$Cs^{5/4}$, since $161/128>5/4$. The complementary term is at most
$C\rho^{m/2}$. On a fixed-band complement a long block gives
$C_{B_*}\rho_{B_*}^{m/2}$.

There is also the independent physical bound at every real frequency,
\[
 |\ell(Q^{m-j}M_sQ^j\mathbf1)|
        \le\nu\{c<s\}=1-\sqrt{1-s^2}\le s^2.
\]
If $|z|\le A+D$ and $|z|\le V$, then
$|z|\le A+\min(D,V)$. Apply this with the Gaussian leading term
as $A$. Its four-frequency integral is $Cs^{5/4}m^{-2}$; the
remainder retains the cap $\min(C_{B_*}\rho_{B_*}^{m/2},Cs^2)$.
The positive interval majorant has Fourier supremum at most
$|J|+2\pi/B_*$ and roof support in $[-B_*,B_*]$. Torus coefficient
extraction and integration over that fixed band therefore give
\eqref{eq:v80-incidence-local}. The cap is essential: an error
independent of $s$ could not be summed against an inverse incidence.
\end{proof}
